\subsubsection{Septic order, \texorpdfstring{$O(\gs^5)$}{O(gs**5)}}
\label{p:sec:hs:septic}

The septic vertex $\vertex{7}$ is the first one that cannot be written with $\Sigma_2$ alone.
This is the order at which the $3$-string field introduced in \cref{p:sec:hs:sextic} becomes indispensable, and at which the pattern repeated at all orders in \cref{p:sec:all-orders} becomes visible.
We first display the structure of the action obtained so far, which allows the auxiliary fields to be eliminated one after the other by Gaussian integrals, exact at tree level.
We then compute the interactions this elimination induces at order $\gs^5$, and fix the cubic coupling of $\vertex{7}$ that cancels them.

\paragraph{Elimination of the auxiliary fields}
\label{p:sec:hs:septic:structure}

\subparagraph{Triangular form of the action.}

The terms of $\caction[\le 6]$ which contain an auxiliary field are, by~\eqref{p:eq:hs5:F} and~\eqref{p:eq:hs6:S6tilde}:
\begin{equation}
  \label{p:eq:hs7:action-6}
  \begin{aligned}
    \caction[\le 6] - S_{\le 3}
      &
      = - \frac{\gs^2}{4!} \, \vertex{4}(\Sigma_2^2)
        + \frac{2 \gs^2}{4!} \, \vertex{4}(\Sigma_2, \Psi^2)
        + \frac{\gs^3}{5!} \, \vertex{5}(\Sigma_2^2, \Psi)
      \\ &
      \quad
      - \frac{\gs^4}{6!} \, \pair[3]{\Sigma_3}{\wprod{3}{3}(\Sigma_3)}
        + \frac{2 \gs^4}{6!} \, \pair[3]{\Sigma_3}{\wprod{3}{3}(\Psi \otimes \Sigma_2)},
  \end{aligned}
\end{equation}
where $S_{\le 3}$ is the kinetic term and the cubic vertex of $\Psi$, common to the cubic action and to the target, and where the sextic term is written through the kinetic operator $\wprod{3}{3}$ of~\eqref{p:eq:hs6:kinetic}.
The rest of the construction relies on two features of this expression.
First, the action is quadratic in each auxiliary field separately: in $\Sigma_3$, with the field-independent operator $\wprod{3}{3}$, and in $\Sigma_2$, with the dressed operator $\kin{2}{\Psi}{3}$ of~\eqref{p:eq:hs5:dressed}.
Second, the source of $\Sigma_3$ is linear in $\Sigma_2$ and $\Psi$.
The source of $\Sigma_2$ involves $\Psi$ only, apart from the coupling to $\Sigma_3$, which is the source of $\Sigma_3$ read the other way.
Once the higher-order auxiliary field $\Sigma_3$ is eliminated, the action is again quadratic in $\Sigma_2$, with a source that depends on $\Psi$ only.
We can therefore eliminate the auxiliary fields one at a time, in decreasing order, $\Sigma_3$ first, and each elimination is a Gaussian integral, exact at tree level.
The value of the action on the solution then comes out in closed form rather than order by order.
We call this the triangular form of the action.


\subparagraph{Eliminating the $3$-string field.}

We eliminated the $3$-string field in~\eqref{p:eq:hs6:eliminate-sigma3}, exactly in $\gs$ and for any $\Sigma_2$.
The solution is $\Sigma_3^{(4)*} = \Psi \otimes \Sigma_2$, and the value of the $\Sigma_3$ terms on it is:
\begin{equation}
  \label{p:eq:hs7:sigma3-value}
    \pair[3]{\wprod{3}{3}(\Psi \otimes \Sigma_2)}{\big( \wprod{3}{3} \big)^{-1} \circ \wprod{3}{3}(\Psi \otimes \Sigma_2)}
    = \dvertex{6}(\Psi \otimes \Sigma_2, \Psi \otimes \Sigma_2),
\end{equation}
which is exactly quadratic in $\Sigma_2$, since the source is linear in $\Sigma_2$.
The inverse of $\wprod{3}{3}$ is taken on the quotient~\eqref{p:eq:hs6:quotient}, under the assumption~\eqref{p:eq:hs6:invertibility-3}.
The action left for $\Sigma_2$ is therefore quadratic too, since the quintic vertex has only two $\Sigma_2$ slots, and the second elimination is again a Gaussian integral.
We expand around the quartic-order solution, $\Sigma_2 = \Psi^2 + \sigma_2^{(2)}$, which does not depend on $\gs$; the solution of this order then follows by completing the square in the fluctuation $\sigma_2^{(2)}$.
The terms of $\caction[\le 6]$ with $\Sigma_3$ eliminated then sort into three groups:
\begin{subequations}
\label{p:eq:hs7:sigma2-groups}
\begin{align}
  \label{p:eq:hs7:sigma2-contact}
  \caction[\le 6] \big|_{\Sigma_3 = \Sigma_3^{(4)*}} - S_{\le 3}
    &
    = \frac{\gs^2}{4!} \, \vertex{4}(\Psi^4)
      + \frac{\gs^3}{5!} \, \vertex{5}(\Psi^5)
      + \frac{\gs^4}{6!} \, \dvertex{6}(\Psi^6)
  \\
  \label{p:eq:hs7:sigma2-quadratic}
  &
  \quad
  - \frac{\gs^2}{4!} \, \pair[2]{\sigma_2^{(2)}}{\kin{2}{\Psi}{4}(\sigma_2^{(2)})}
  \\
  \label{p:eq:hs7:sigma2-linear}
  &
  \quad
  + \frac{\gs^3}{60} \, \Big[
    \vertex{5}(\sigma_2^{(2)}, \Psi^2, \Psi)
    + \frac{\gs}{6} \, \dvertex{6}(\Psi^3, \Psi \otimes \sigma_2^{(2)})
  \Big],
\end{align}
\end{subequations}
where the first line collects the contact terms, the second is the quadratic form of the fluctuation and the third is linear in it.
The contact terms are the target through order $\gs^4$ except for the sextic deficit: $\dvertex{6}(\Psi^6)$ differs from $\vertex{6}(\Psi^6)$ by $- \tfrac{6}{5} \vertex{5 \cdot 5}(\Psi^6)$, which the exchange term of the quintic order compensates, as in~\eqref{p:eq:hs6:eliminate}.
The quadratic form is that of the dressed kinetic operator $\kin{2}{\Psi}{4}$ of~\eqref{p:eq:hs6:dressed-2}.
It is the quartic form dressed by the quintic vertex, as in~\eqref{p:eq:hs5:dressed}, and by the exchange of $\Sigma_3$ between two sextic couplings: by~\eqref{p:eq:hs7:sigma3-value}, this exchange is the last term of $\kin{2}{\Psi}{4}$.
The linear term defines the source of the fluctuation, the $2$-string state $\beta_\Psi \in \Hn{2}$ such that
\begin{equation}
  \label{p:eq:hs7:source-2}
  \pair[2]{\beta_\Psi}{A}
    := \vertex{5}(A, \Psi^2, \Psi) + \frac{\gs}{6} \, \dvertex{6}(\Psi^3, \Psi \otimes A),
  \qquad
  \beta_\Psi
    = \sprod{3}{2}(\Psi^3) + O(\gs),
\end{equation}
so that~\eqref{p:eq:hs7:sigma2-linear} reads $\tfrac{\gs^3}{60} \pair[2]{\beta_\Psi}{\sigma_2^{(2)}}$.
Completing the square, we get the value of $\caction[\le 6]$ on the joint solution as a single expression, exact in $\gs$:
\begin{equation}
  \label{p:eq:hs7:closed-form}
  \caction[\le 6] \big|_{\Sigma_3 = \Sigma_3^{(4)*}, \, \Sigma_2 = \Sigma_2^{(4)*}} - S_{\le 6}
    = \frac{\gs^4}{600} \Big[ \pair[2]{\beta_\Psi}{\big( \kin{2}{\Psi}{4} \big)^{-1}(\beta_\Psi)} - \vertex{5 \cdot 5}(\Psi^6) \Big],
\end{equation}
where the last term is the sextic deficit of~\eqref{p:eq:hs7:sigma2-contact}.
The right-hand side is a single function of $\gs$, rational at any finite level truncation: it describes the the whole sequence of interactions induced by the two auxiliary fields.
It is as explicit as the inverse of the dressed operator, which is a Neumann series for small $\gs$ as in \cref{p:sec:hs:quintic}.
It holds at finite $\gs$ when $\wprod{3}{3}$ and $\kin{2}{\Psi}{4}$ are invertible on their quotients.
We checked~\eqref{p:eq:hs7:closed-form} against a direct solution of the joint equations of motion in the finite-dimensional models of \cref{p:app:finite-models}, with generic vertices.

\subparagraph{Induced interactions at order $\gs^5$.}

We expand~\eqref{p:eq:hs7:closed-form} to first order in $\gs$ beyond the sextic order, with the inverse of the dressed operator:
\begin{equation}
\big( \kin{2}{\Psi}{4} \big)^{-1}
  = \sprod{2}{2}^{-1}
    + \frac{\gs}{5} \, \sprod{2}{2}^{-1} \circ \sprod{3}{2}(\Psi, \cdot) \circ \sprod{2}{2}^{-1}
    + O(\gs^2).
\end{equation}
The $\Sigma_3$ exchange inside $\kin{2}{\Psi}{4}$ enters only at order $\gs^6$.
We get:
\begin{equation}
  \label{p:eq:hs7:induced}
  \caction[\le 6] \big|_{\Sigma_3 = \Sigma_3^{(4)*}, \, \Sigma_2 = \Sigma_2^{(4)*}}
    = S_{\le 6}
      + \frac{\gs^5}{600} \Big[ \frac{1}{3} \, \mathcal S_2(\Psi^7) - \frac{1}{5} \, \mathcal S_1(\Psi^7) \Big]
      + O(\gs^6),
\end{equation}
with the two seven-point structures
\begin{subequations}
  \label{p:eq:hs7:structures}
  \begin{align}
    \label{p:eq:hs7:S1}
    \mathcal S_1(\Psi^7)
      &
      := \vertex{5}\Big( \sprod{2}{2}^{-1} \circ \sprod{3}{2}(\Psi^3), \, \sprod{2}{2}^{-1} \circ \sprod{3}{2}(\Psi^3), \, \Psi \Big),
    \\
    \label{p:eq:hs7:S2}
    \mathcal S_2(\Psi^7)
      &
      := \vertex{6}\Big( \Psi^4, \, \sprod{2}{2}^{-1} \circ \sprod{3}{2}(\Psi^3) \Big).
  \end{align}
\end{subequations}
Both are exchange diagrams of $\Sigma_2$ between interaction vertices, and neither collapses, by the criterion of \cref{p:sec:hs:sextic}.
The first is a chain of three quintic vertices joined by two $\Sigma_2$ lines, the two end vertices carrying three external $\Psi$ each and the middle one a single $\Psi$.
It is the next term of the Neumann series~\eqref{p:eq:hs5:saddle-series} inserted in the exchange term of~\eqref{p:eq:hs5:identity}, which gives $+ \tfrac{\gs^5}{3000} \mathcal S_1$.
The second is a sextic vertex with four external $\Psi$ joined by a $\Sigma_2$ line to a quintic vertex with three.
It comes from the reduced sextic vertex evaluated on the quintic-order solution, that is, from the $\Sigma_3$ exchange~\eqref{p:eq:hs7:sigma3-value} with one $\Sigma_2$ slot of $\dvertex{6}(\Psi \otimes \Sigma_2, \Psi \otimes \Sigma_2)$ receiving the correction $\tfrac{\gs}{5} \sprod{2}{2}^{-1} \circ \sprod{3}{2}(\Psi^3)$ of~\eqref{p:eq:hs5:saddle-series}.
This cross term is $\tfrac{2 \gs^5}{5 \cdot 6!} \dvertex{6}(\Psi^3, \Psi \otimes B)$ with $B := \sprod{2}{2}^{-1} \circ \sprod{3}{2}(\Psi^3)$.
By $\dvertex{6}(\Psi^3, \Psi \otimes B) = \vertex{6}(\Psi^4, B) - \tfrac{6}{5} \vertex{5 \cdot 5}(\Psi^3, \Psi \otimes B)$, whose exchange part is $\mathcal S_1$ again, it equals $\tfrac{\gs^5}{1800} [\mathcal S_2 - \tfrac{6}{5} \mathcal S_1]$.
The two contributions add up to~\eqref{p:eq:hs7:induced}.
No other term enters at order $\gs^5$: the correction of order $\gs^2$ to the solution for $\Sigma_2$ from the sextic term contributes only at order $\gs^6$, by stationarity of $\caction[\le 5]$ at its own solution, as at the sextic order.
In both structures, each triple of external legs is attached to one and the same vertex.
  We use this property to build the reduced vertex below.

\paragraph{Septic coupling}
\label{p:sec:hs:septic:dressed}

\subparagraph{Reduced septic vertex.}


The septic order must reproduce $\vertex{7}$ and cancel the induced term of~\eqref{p:eq:hs7:induced} at the same time.
A term added at order $\gs^5$ contributes to the eliminated action at that order through its value on the diagonal $\Sigma_k = \Psi^k$, where it must be of degree seven in $\Psi$, and it must be cubic in the fields.
The only partitions of seven strands over three fields drawn from $\{\Psi, \Sigma_2, \Sigma_3\}$ are $7 = 3 + 3 + 1$ and $7 = 3 + 2 + 2$.
We take the first partition, in which the $3$-string field plays the role $\Sigma_2$ played in the quintic coupling:
\begin{equation}
  \label{p:eq:hs7:S7tilde}
  \caction[7]
    := \frac{\gs^5}{7!} \, \dvertex{7}(\Sigma_3, \Sigma_3, \Psi),
  \qquad
  \caction[\le 7]
    := \caction[\le 6] + \caction[7],
\end{equation}
with a trilinear form $\dvertex{7}$ on $\Hn{3} \times \Hn{3} \times \Hs$, symmetric in its two $3$-string arguments, which the matching determines.
  We discuss the second partition below.
We define the reduced septic vertex, for $A, B \in \Hn{3}$ and $C \in \Hs$:
\begin{equation}
  \label{p:eq:hs7:dressed-vertex}
  \begin{aligned}
    \dvertex{7}(A, B, C)
      &
      := \vertex{7}(A, B, C)
        + \frac{42}{25} \, \vertex{5}\Big( \sprod{2}{2}^{-1} \circ \sprod{3}{2}(A), \, \sprod{2}{2}^{-1} \circ \sprod{3}{2}(B), \, C \Big)
      \\ &
      \quad
      - \frac{14}{5} \cdot \frac{1}{2} \Big[
        \vertex{6}\Big( A, C, \sprod{2}{2}^{-1} \circ \sprod{3}{2}(B) \Big)
        + \vertex{6}\Big( B, C, \sprod{2}{2}^{-1} \circ \sprod{3}{2}(A) \Big)
      \Big],
  \end{aligned}
\end{equation}
where $\sprod{3}{2}(A)$ fills the three input slots of $\sprod{3}{2}$ with the $3$-string state $A$, and where $\vertex{6}(A, C, D)$ fills three slots with $A$, one with $C$ and two with $D \in \Hn{2}$.
The second term is the chain of three quintic vertices of $\mathcal S_1$, the two end vertices read with their three single-string slots as one $\Sigma_3$ slot.
The third term is the sextic--quintic exchange of $\mathcal S_2$, with one $\Sigma_3$ in the sextic vertex and the other in the quintic one, symmetrized.
On the diagonal
\begin{equation}
  \dvertex{7}(\Psi^3, \Psi^3, \Psi)
    = \vertex{7}(\Psi^7)
      + \frac{42}{25} \, \mathcal S_1(\Psi^7)
      - \frac{14}{5} \, \mathcal S_2(\Psi^7),
\end{equation}
by~\eqref{p:eq:hs7:structures}.
\begin{figure}[htp]
  \centering
  \begin{tikzpicture}
    \begin{feynman}
      \vertex (a1) at (0, 0) {$\Sigma_3$};
      \vertex (a) at (1.2, 0) [dot];
      \vertex (b) at (2.8, 0) [dot];
      \vertex (c) at (4.4, 0) [dot];
      \vertex (c1) at (5.6, 0) {$\Sigma_3$};
      \vertex (p) at (2.8, 1.2) {$\Psi$};
      \diagram* {
      (a1) -- [plain, double distance=2.5pt] (a) -- [plain, double, edge label'={$\Sigma_2$}] (b) -- [plain, double, edge label'={$\Sigma_2$}] (c) -- [plain, double distance=2.5pt] (c1),
      (b) -- [plain] (p),
      };
      \vertex (d1) at (7.5, 0) {$\Sigma_3$};
      \vertex (d) at (8.7, 0) [dot];
      \vertex (e) at (10.3, 0) [dot];
      \vertex (e1) at (11.5, 0) {$\Sigma_3$};
      \vertex (q) at (8.7, 1.2) {$\Psi$};
      \diagram* {
      (d1) -- [plain, double distance=2.5pt] (d) -- [plain, double, edge label'={$\Sigma_2$}] (e) -- [plain, double distance=2.5pt] (e1),
      (d) -- [plain] (q),
      };
    \end{feynman}
  \end{tikzpicture}
  \caption{The two exchange terms of the reduced septic vertex~\eqref{p:eq:hs7:dressed-vertex}: a thick double line is a $\Sigma_3$ leg (three strands), a double line a $\Sigma_2$ line carrying $\sprod{2}{2}^{-1}$. Left, three quintic vertices on two $\Sigma_2$ lines; right, a sextic vertex and a quintic vertex on one $\Sigma_2$ line. In both, the three strands of each $\Sigma_3$ leg enter one and the same vertex.}
  \label{p:fig:hs7:dressed-vertex}
\end{figure}

The elimination proceeds as before, $\Sigma_3$ first.
The septic coupling is quadratic in $\Sigma_3$ and does not change its source, so the $3$-string sector stays Gaussian, with the dressed kinetic operator:
\begin{equation}
  \label{p:eq:hs7:dressed-3}
  \pair[3]{A}{\kin{3}{\Psi}{5}(B)}
    := \pair[3]{A}{\wprod{3}{3}(B)} - \frac{\gs}{7} \, \dvertex{7}(A, B, \Psi),
  \qquad
  A, B \in \Hn{3}.
\end{equation}
The equation of motion of $\Sigma_3$ and its solution are:
\begin{equation}
  \label{p:eq:hs7:saddle-3}
  \begin{aligned}
    \mathcal E^{(5)}_{\Sigma_3}
      &
      := \wprod{3}{3}(\Psi \otimes \Sigma_2) - \kin{3}{\Psi}{5}(\Sigma_3)
      = 0,
    \\
    \Sigma_3^{(5)*}
      &
      := \big( \kin{3}{\Psi}{5} \big)^{-1} \circ \wprod{3}{3}(\Psi \otimes \Sigma_2),
  \end{aligned}
\end{equation}
the inverse existing as a Neumann series in $\tfrac{\gs}{7} (\wprod{3}{3})^{-1} \circ \dvertex{7}(\, \cdot \,, \cdot \,, \Psi)$ for small $\gs$.
The solution is now $\Psi \otimes \Sigma_2$ plus a correction of order $\gs$, as $\Sigma_2^{(3)*}$ at the quintic order.
The value of the $\Sigma_3$ terms on the solution is the completed square:
\begin{equation}
  \label{p:eq:hs7:sigma3-value-5}
  \begin{aligned}
    &
    \frac{\gs^4}{6!} \, \pair[3]{\wprod{3}{3}(\Psi \otimes \Sigma_2)}{\big( \kin{3}{\Psi}{5} \big)^{-1} \circ \wprod{3}{3}(\Psi \otimes \Sigma_2)}
    \\ &
    \qquad
    = \frac{\gs^4}{6!} \, \dvertex{6}(\Psi \otimes \Sigma_2, \Psi \otimes \Sigma_2)
      + \frac{\gs^5}{7!} \, \dvertex{7}(\Psi \otimes \Sigma_2, \Psi \otimes \Sigma_2, \Psi)
      + O(\gs^6),
  \end{aligned}
\end{equation}
where the first term of the Neumann series gives back~\eqref{p:eq:hs7:sigma3-value} and the second is the septic coupling.
The order-$\gs^5$ contribution of the septic coupling on-shell is thus its naive contact value, as for the quintic coupling in~\eqref{p:eq:hs5:constant}.
The action left for $\Sigma_2$ is still exactly quadratic, since $\Sigma_3^{(5)*}$ is linear in $\Sigma_2$.
Its solution $\Sigma_2^{(5)*}$ differs from $\Sigma_2^{(4)*}$ at order $\gs^3$.
By stationarity this does not affect order $\gs^5$, and at that order the septic term enters only through its value at $\Sigma_2 = \Psi^2$.

The target at this order is:
  \begin{equation}
  S_{\le 7}
    = S_{\le 6} + \frac{\gs^5}{7!} \, \vertex{7}(\Psi^7).
\end{equation}
  When we add the contact value of the septic coupling to~\eqref{p:eq:hs7:induced}, the matching with this target requires the single condition:
\begin{equation}
  \label{p:eq:hs7:coefficient}
  \frac{\gs^5}{600} \Big[ \frac{1}{3} \, \mathcal S_2 - \frac{1}{5} \, \mathcal S_1 \Big]
    + \frac{\gs^5}{7!} \, \dvertex{7}(\Psi^7)
    = \frac{\gs^5}{7!} \, \vertex{7}(\Psi^7),
\end{equation}
that is,
\begin{equation}
  \label{p:eq:hs7:diagonal}
  \dvertex{7}(\Psi^7)
    = \vertex{7}(\Psi^7) + \frac{42}{25} \, \mathcal S_1(\Psi^7) - \frac{14}{5} \, \mathcal S_2(\Psi^7),
\end{equation}
which is the diagonal value of~\eqref{p:eq:hs7:dressed-vertex}.
The matching determines the three coefficients $(1, \tfrac{42}{25}, - \tfrac{14}{5})$, since $\vertex{7}(\Psi^7)$, $\mathcal S_1$ and $\mathcal S_2$ are independent functions of $\Psi$ in a finite-dimensional model with generic vertices.
We assume this independence for the string vertices, as at the sextic order.
  Changing one coefficient then leaves a seven-point term proportional to the corresponding invariant.
With these coefficients, we get:
\begin{equation}
  \label{p:eq:hs7:eliminate}
  \caction[\le 7] \big|_{\Sigma_3 = \Sigma_3^{(5)*}, \, \Sigma_2 = \Sigma_2^{(5)*}}
    = S_{\le 7} + O(\gs^6),
\end{equation}
and the septic order closes: the exchange diagrams induced at order $\gs^5$ are canceled by the exchange parts of $\dvertex{7}$, one structure at a time, and the septic vertex is reproduced.
We checked~\eqref{p:eq:hs7:eliminate} in the finite-dimensional models, where the residual scales as $\gs^6$ and the matching through order $\gs^4$ is untouched.
The remainder of order $\gs^6$ is the next induced interaction, which the octic order absorbs in its turn (\cref{p:sec:hs:octic}).

\subparagraph{Dressed kinetic operators.}

The septic coupling dresses the kinetic operator of $\Sigma_3$ with $\Psi$, through~\eqref{p:eq:hs7:dressed-3}, as the quintic coupling did for $\Sigma_2$.
We can read its two exchange parts in terms of the exchanges of the lower orders, and this reading shows the pattern of the higher orders.
The chain term of~\eqref{p:eq:hs7:dressed-vertex} enters $\kin{3}{\Psi}{5}$ with the coefficient $- \tfrac{6}{5} \cdot \tfrac{\gs}{5}$.
This is the order-$\gs$ term of the exchange part of $\wprod{3}{3}$ with the bare $\Sigma_2$ propagator replaced by the dressed one,
\begin{equation}
  \label{p:eq:hs7:running}
  \begin{aligned}
    \sprod{2}{3} \circ \big( \kin{2}{\Psi}{3} \big)^{-1} \circ \sprod{3}{2}
      &
      = \sprod{2}{3} \circ \sprod{2}{2}^{-1} \circ \sprod{3}{2}
        + \frac{\gs}{5} \, \sprod{2}{3} \circ \sprod{2}{2}^{-1} \circ \sprod{3}{2}(\Psi, \cdot) \circ \sprod{2}{2}^{-1} \circ \sprod{3}{2}
        \\ & \qquad
        + O(\gs^2),
  \end{aligned}
\end{equation}
by the Neumann series~\eqref{p:eq:hs5:neumann}.
At this order, the chain term of the dressing is thus the first Taylor coefficient of the operator obtained from $\wprod{3}{3}$ by replacing the bare $\Sigma_2$ propagator with the dressed one, and $\wprod{3}{3}$ is the value of that operator at $\gs = 0$.
The higher coefficients are not tested here.
The sextic--quintic term, whose coefficient $+ \tfrac{2 \gs}{5}$ follows from the matching, is the $\Sigma_2$ exchange between the contact part $\vertex{6}$ of the sextic source vertex of $\Sigma_3$ and the quintic vertex read as $\sprod{3}{2}(\Sigma_3)$.
Conversely, the dressed operator of $\Sigma_2$ at this order is read off the action once $\Sigma_3$ is eliminated with~\eqref{p:eq:hs7:sigma3-value-5}:
\begin{equation}
  \label{p:eq:hs7:dressed-2}
  \begin{aligned}
    \pair[2]{A}{\kin{2}{\Psi}{5}(B)}
      &
      := \vertex{4}(A, B)
        - \frac{\gs}{5} \, \vertex{5}(A, B, \Psi)
      \\ &
      \quad
      - \frac{\gs^2}{30} \, \pair[3]{\wprod{3}{3}(\Psi \otimes A)}{\big( \kin{3}{\Psi}{5} \big)^{-1} \circ \wprod{3}{3}(\Psi \otimes B)},
  \end{aligned}
\end{equation}
This is the operator $\kin{2}{\Psi}{4}$ of~\eqref{p:eq:hs6:dressed-2} with the $\Sigma_3$ propagator dressed in its turn.
Each order thus dresses the kinetic operators of the lower-order auxiliary fields by the exchanges of the higher-order ones, and the propagators inside these exchanges are themselves the dressed ones.
The diagonal value of every reduced vertex is fixed in the same way: it is the vertex of the target minus what the previous orders induce:
\begin{equation}
  \label{p:eq:hs7:deficit}
  \dvertex{n}(\Psi^n)
    = \vertex{n}(\Psi^n) - n! \, \big[ \gs^{n - 2} \big] \, \caction[\le n - 1] \big|_{\text{solution}},
\end{equation}
where the cubic action of the previous order is evaluated on the solution of its auxiliary fields and $[\gs^{n - 2}]$ extracts the coefficient of $\gs^{n - 2}$; the reduced sextic vertex~\eqref{p:eq:hs6:dressed-vertex} and the condition~\eqref{p:eq:hs7:coefficient} are its cases $n = 6$ and $n = 7$.

\paragraph{Choices in the construction}
\label{p:sec:hs:septic:alternatives}

\subparagraph{Choice of the partition.}

The matching~\eqref{p:eq:hs7:coefficient} constrains the septic term only through its value on the diagonal, a septic polynomial in $\Psi$.
Any trilinear form with this diagonal value closes the order, and two forms differing by one that vanishes on the diagonal give the same eliminated action at this order.
The exchange parts of~\eqref{p:eq:hs7:dressed-vertex} are unique only up to this freedom, as those of the reduced sextic vertex were, and the freedom now includes the partition of the strands.
The second partition, $7 = 3 + 2 + 2$, gives a coupling $\vertex{7}(\Sigma_3, \Sigma_2, \Sigma_2)$ which, reduced by the same two structures with the same diagonal value, closes the order as well by the argument just given.
So does any combination of the two partitions whose coefficients add up to one.
To realize $\mathcal S_1$ and $\mathcal S_2$ in the second partition, the two strands of a $\Sigma_2$ slot must be split between two constituent vertices.
We prefer the first partition for three reasons.
In the first partition all the strands of each slot sit inside one constituent vertex, as in \cref{p:fig:hs7:dressed-vertex}, so that each exchange part of the reduced vertex is one of the induced diagrams~\eqref{p:eq:hs7:structures} with its triples of legs read as $\Sigma_3$ slots, off the diagonal as well as on it.
The first partition also preserves the triangular form.
A coupling $\vertex{7}(\Sigma_3, \Sigma_2, \Sigma_2)$ makes the source of $\Sigma_3$ quadratic in $\Sigma_2$, the action left after eliminating $\Sigma_3$ becomes quartic in $\Sigma_2$, and the closed form~\eqref{p:eq:hs7:closed-form} is lost at the next order.
With the first partition, finally, $\Sigma_3$ comes with an even and an odd coupling, $\dvertex{6}$ and $\dvertex{7}$, as $\Sigma_2$ came with $\vertex{4}$ and $\vertex{5}$.
We therefore keep the first partition,~\eqref{p:eq:hs7:S7tilde} with the reduced vertex~\eqref{p:eq:hs7:dressed-vertex}.

\subparagraph{The necessity of the $3$-string field.}

We could ask whether the $3$-string field could be introduced at this order, where it is first indispensable, rather than at the sextic order.
Suppose that the sextic order had been closed with $\Sigma_2$ alone, with a sextic coupling in $\Sigma_2^3$, see \cref{p:sec:hs:sextic}, and that $\Sigma_3$ were introduced here with a kinetic term and a single coupling, $\vertex{7}(\Sigma_3, \Sigma_2, \Sigma_2)$ or $\vertex{7}(\Sigma_3, \Sigma_3, \Psi)$.
The only kinetic term for $\Sigma_3$ built from the vertices is of sextic type, $\vertex{6}(\Sigma_3^2)$ or its reduced form, and is of order $\gs^4$.
Without a source at that order, the solution for $\Sigma_3$ is of order $\gs$ and the value of the coupling on it is of order $\gs^6$ at least.
The contact term $\vertex{7}(\Psi^7)$ is then never produced at order $\gs^5$, for any kernel and any coefficient.
Within cubic actions in $\Psi$, $\Sigma_2$ and $\Sigma_3$ with kinetic terms of the same sign, a kinetic term and a source for $\Sigma_3$ at the sextic order, giving it a solution of order one, are therefore required before the septic order can close.
The completed square~\eqref{p:eq:hs6:DeltaS6} is the form we use.
Within this class, a sextic order closed with $\Sigma_2$ alone has to be rebuilt with $\Sigma_3$ as soon as $\vertex{7}$ is reached: that form reproduces the same eliminated action through order $\gs^4$ and does not extend beyond it, the distinction made in \cref{p:sec:hs:sextic}.

This is why we introduce each auxiliary field through a completed square at the even order $\gs^{2 k - 2}$ and let the odd vertex dress its kinetic operator at the next order.
The pattern is established here for $k = 3$ and repeated in \cref{p:sec:all-orders}.
The septic order thus introduces no new field: it couples $\Sigma_3$ beyond the completed square of the sextic order, and gives it its first dressed kinetic operator.

\subparagraph{Kernel condition.}

The septic coupling must be a function on the quotient~\eqref{p:eq:hs6:quotient} on which $\Sigma_3$ lives, as the quintic coupling had to be a function on $\Hn{2} / \ker \sprod{2}{2}$.
The sextic source $\wprod{3}{3}(\Psi \otimes \Sigma_2)$ pairs to zero with every kernel element by the symmetry of $\wprod{3}{3}$, so it descends automatically, whereas the septic coupling does not.
From this order on we make a further standing assumption, besides those of \cref{p:sec:hs:quintic,p:sec:hs:sextic}.
The kernel of the kinetic operator with which each auxiliary field is introduced, $\wprod{3}{3}$ for $\Sigma_3$ here and its analogue for $\Sigma_k$ in \cref{p:sec:all-orders}, consists of joint null vectors of every vertex in which that field appears.
It need only hold for one choice of the terms vanishing on the diagonal, up to which the reduced vertices are defined.
At this order the condition reads:
\begin{equation}
  \label{p:eq:hs7:kernel-condition}
  \forall N_3 \in \ker \wprod{3}{3}, \;
  \forall A \in \Hn{3}, \;
  \forall B \in \Hs:
  \quad
  \dvertex{7}(N_3, A, B)
    = 0,
\end{equation}
and it makes $\dvertex{7}(\, \cdot \,, \cdot \,, \Psi)$, hence $\kin{3}{\Psi}{5}$ and its inverse, descend to the quotient.
It holds trivially when $\wprod{3}{3}$ has no kernel, and it is a condition on the choice of the reduced vertex rather than on the vertices alone, since $\dvertex{7}$ may be changed by terms vanishing on the diagonal.
When it fails for some kernel direction $N_3$, that direction has no kinetic term at order $\gs^4$ and acquires a quadratic term from the background at order $\gs^5$, $\tfrac{\gs^5}{7!} \dvertex{7}(N_3, N_3, \Psi)$.
Integrating it out divides by this field-dependent mass, and the eliminated action acquires a term rational in $\Psi$, $- \tfrac{\gs^5}{7!} \dvertex{7}(N_3, \Psi^3, \Psi)^2 / \dvertex{7}(N_3, N_3, \Psi)$ when the denominator does not vanish, which no polynomial vertex can match.
We established this in a finite-dimensional model.
Under the assumption, the auxiliary fields reproduce a polynomial action order by order.
When it fails along a direction where the denominator does not vanish, the eliminated action is not polynomial.
We do not try to settle here whether the string vertices satisfy this condition.
\Cref{p:sec:all-orders} and \cref{p:sec:gauge} refer to it in this form.

\paragraph{Gauge transformation}
\label{p:sec:hs:septic:gauge}

\subparagraph{Rotation of the fluctuations.}

With the septic coupling, the action decomposes exactly, off-shell, into the target, the induced interactions and two squares:
\begin{equation}
  \label{p:eq:hs7:identity}
  \begin{aligned}
    \caction[\le 7]
      &
      = S_{\le 7}
        + \big( \text{exchange terms of order } \gs^6 \big)
      \\ &
      \quad
      - \frac{\gs^4}{6!} \, \pair[3]{\sigma_3^{(5)}}{\kin{3}{\Psi}{5}(\sigma_3^{(5)})}
      - \frac{\gs^2}{4!} \, \pair[2]{\sigma_2^{(5)}}{\kin{2}{\Psi}{5}(\sigma_2^{(5)})},
  \end{aligned}
\end{equation}
where the exchange terms of order $\gs^6$ are the induced interaction of~\eqref{p:eq:hs7:eliminate}, a function of $\Psi$ alone.
The fluctuations are:
\begin{equation}
  \label{p:eq:hs7:deviations}
  \sigma_3^{(5)}
    := \Sigma_3 - \Sigma_3^{(5)*},
  \qquad
  \sigma_2^{(5)}
    := \Sigma_2 - \Sigma_2^{(5)*},
\end{equation}
the first from the solution~\eqref{p:eq:hs7:saddle-3} at the given $\Sigma_2$, the second from the solution of the $\Sigma_2$ equation once $\Sigma_3$ is eliminated.
Both squares come with dressed kinetic operators.
As at the quintic order, we fix the gauge law by demanding that each square be exactly invariant.
Since the action is Gaussian in both fluctuations, the transformation laws are explicit rotations:
\begin{equation}
  \label{p:eq:hs7:gauge-deviations}
  \begin{aligned}
    \delta_\Lambda \sigma_3^{(5)}
      &
      = - \frac{1}{2} \, \big( \kin{3}{\Psi}{5} \big)^{-1} \circ \big( \delta_\Lambda \kin{3}{\Psi}{5} \big) (\sigma_3^{(5)})
      = \frac{\gs}{14} \, \big( \kin{3}{\Psi}{5} \big)^{-1} \circ \dvertex{7}(\, \cdot \,, \sigma_3^{(5)}, \delta_\Lambda \Psi),
    \\
    \delta_\Lambda \sigma_2^{(5)}
      &
      = - \frac{1}{2} \, \big( \kin{2}{\Psi}{5} \big)^{-1} \circ \big( \delta_\Lambda \kin{2}{\Psi}{5} \big) (\sigma_2^{(5)}),
  \end{aligned}
\end{equation}
where $\dvertex{7}(\, \cdot \,, \sigma_3^{(5)}, \delta_\Lambda \Psi)$ denotes the $3$-string state paired with the open slot, as in~\eqref{p:eq:hs6:eom}.
Each auxiliary field transforms as its solution plus the rotation of its fluctuation:
\begin{equation}
  \label{p:eq:hs7:gauge-fields}
  \delta_\Lambda \Sigma_3
    = \delta_\Lambda \Sigma_3^{(5)*} + \delta_\Lambda \sigma_3^{(5)},
  \qquad
  \delta_\Lambda \Sigma_2
    = \delta_\Lambda \Sigma_2^{(5)*} + \delta_\Lambda \sigma_2^{(5)}.
\end{equation}
The variation of $\kin{2}{\Psi}{5}$ collects those of the quintic term, of the sextic source $\wprod{3}{3}(\Psi \otimes \cdot)$ and of the inverse of $\kin{3}{\Psi}{5}$ in~\eqref{p:eq:hs7:dressed-2}.
With these transformation laws, varying~\eqref{p:eq:hs7:identity} gives:
\begin{equation}
  \label{p:eq:hs7:delta-S-equal}
  \delta_\Lambda \caction[\le 7] - \delta_\Lambda \big[ S_{\le 7} + \big( \text{exchange terms of order } \gs^6 \big) \big]
    = 0
\end{equation}
identically, off-shell, for any variation $\delta_\Lambda \Psi$ and exactly in $\gs$.
Since $\delta_\Lambda S_{\le 7} = O(\gs^6)$, the cubic action is invariant up to order $\gs^6$, the order of the truncation it reproduces.
We checked the identity, and the residuals left by dropping either rotation, in the finite-dimensional model.
At the sextic order, the fluctuation of $\Sigma_3$ needed no rotation, since its kinetic operator was independent of $\Psi$.
From this order on, both fluctuations rotate.
The rotations are homogeneous, $\delta_\Lambda X \propto X$, and vanish on-shell.
Their algebra and the Batalin--Vilkovisky structure of the cubic action with both auxiliary fields are worked out in \cref{p:sec:gauge} through the sextic order.
A flat representative for the square of $\sigma_3^{(5)}$, expected as for $\Sigma_2$, has not been constructed.

\subparagraph{Reality.}

The fields are subject to the conditions $\Sigma_3^\star = - \Sigma_3$ and $\Sigma_2^\star = \Sigma_2$ of \cref{p:app:conventions}, and the septic coupling is compatible with them for the reason given at the quintic and sextic orders.
The reduced vertex $\dvertex{7}$ is built from $\vertex{5}$, $\vertex{6}$, $\vertex{7}$, $\sprod{3}{2}$ and $\sprod{2}{2}^{-1}$ with real coefficients, so the septic term is real on real fields by the reality of the vertices~\eqref{p:eq:hs4:vertex-real}, and $\kin{3}{\Psi}{5}$ is a real form when $\Psi$ is real.
The map $(\kin{3}{\Psi}{5})^{-1} \circ \wprod{3}{3}$ then commutes with $\star$, by conjugating the relation $\pair[3]{A}{\kin{3}{\Psi}{5}(B')} = \dvertex{6}(A, B)$, with $B' := (\kin{3}{\Psi}{5})^{-1} \circ \wprod{3}{3}(B)$, that characterizes it, as at the quintic order.
Hence the solution $\Sigma_3^{(5)*}$ is real when $\Psi$ and $\Sigma_2$ are, and the rotations~\eqref{p:eq:hs7:gauge-deviations} map the real locus to itself.
The argument at every order is given in \cref{p:app:reality}.
